\chapter*{Kata Pengantar}
Buku ini merupakan catatan pribadi saya yang kemudian saya buat menjadi publik.
Siapa tahu catatan ini bermanfaat bagi banyak orang.
Ini saya lakukan karena melihat banyak orang yang belum paham
tentang ``Artificial Intelligence'' atau AI (tetapi banyak bicara tentang itu).
Atau bahkan ada banyak orang yang kebingungan dengan AI.
Mudah-mudahan buku singkat ini dapat mencerahkan.

Buku ini bukan untuk buku teks, meskipun dapat juga digunakan sebagai
bagian dari pengantar AI.

Dalam buku ini saya akan menggunakan istilah ``kecerdasan buatan''
sebagai terjemahan dari ``Artificial Intelligence''.
Tentu saja ada terjemahan yang lain dan ini masih bisa kita diskusikan.
Hanya saja karena ini catatan pribadi saya, maka istilah inilah yang
akan saya gunakan dalam buku ini.

\section*{Tentang Penulis Buku Ini}
Saya sendiri sempat belajar AI di tahun 1988.
Kalau itu dia masih dikenal dengan nama ``(Artificial) Neural Networks''.
Selain itu saya juga belajar ``formal logic'' yang sebetulnya juga merupakan
salah satu ``madzhab'' dari AI. (Tentang ini akan kita bahas panjang lebar.)
Demikian pula saya juga belajar statistik yang menjadi landasan dari
``Bayesian'' AI. Intinya, saya punya latar belakang AI yang cukup
teknis dan cukup mendalam.
Selain itu saya juga membuat beberapa perusahaan AI.
Jadi dapat dikatakan saya tidak hanya sekedar berteori, tetapi juga
membuat implementasi praktis dari AI.

Semoga bermanfaat.

Bandung 2025

Budi Rahardjo @rahard

